\chapter{Применение методов машинного обучения в управлении малым и средним бизнесом: обзор литературы}

\section{Данные как управленческий ресурс в малом и среднем бизнесе}

Малые и средние предприятия (МСП) исторически находятся в невыгодном положении относительно крупного бизнеса с точки зрения доступа к аналитическим инструментам. По данным ОЭСР~\cite{oecd2021sme}, более 65\% МСП не используют цифровые инструменты обработки данных на регулярной основе, а управленческие решения принимаются преимущественно на основе интуиции и разрозненных отчётов. Между тем объём транзакционных данных, накапливаемых даже небольшим производственным предприятием, вполне достаточен для построения базовых прогнозных моделей: ежедневные банковские операции, отчёты об отгрузках, данные о возвратах и платёжный календарь образуют структурированный временно́й ряд за 12--24 месяца.

Разрыв между наличием данных и их практическим применением объясняется не столько технической сложностью ML-методов, сколько отсутствием воспроизводимых аналитических конвейеров, адаптированных к ресурсным ограничениям МСП~\cite{sculley2015hidden}. Крупный бизнес располагает выделенными командами data science и корпоративными хранилищами данных; МСП работает с выгрузками из учётной системы, банковскими выписками в Excel и ручными сводными таблицами. Практическая ценность данного различия в том, что ML-решение для МСП должно быть по определению простым, интерпретируемым и устойчивым к нерегулярности данных --- принципиально иные требования, чем у промышленных ML-систем крупного ритейла~\cite{makridakis2022m5}.

Роль данных в операционном управлении МСП хорошо описана в работах по управленческому учёту~\cite{whittington2012principles}: ключевым инструментом является cash flow statement, поскольку он фиксирует реальное движение денег, а не бухгалтерские начисления. Для производственных предприятий пищевой промышленности, где маржа традиционно невысока, а переговорная позиция по отношению к крупным ритейлерам слабая~\cite{sme_profitability_ru}, прозрачность денежного потока становится условием выживания, а не конкурентным преимуществом.

\section{Автоматическая классификация финансовых транзакций}

Назначения платежей в банковских выписках представляют собой неструктурированный текст со специфической лексикой: аббревиатуры, номера договоров, ИНН контрагентов, произвольные формулировки. Семантическая разметка таких текстов --- задача, находящаяся на пересечении автоматической категоризации текстов~\cite{sebastiani2002machine} и финансовой аналитики.

Классическая работа Себастьяни~\cite{sebastiani2002machine} систематизирует подходы к автоматической категоризации текстов в диапазоне от логических правил (rule-based) до вероятностных методов и классификаторов на основе опорных векторов. В контексте банковских транзакций rule-based подход остаётся доминирующим благодаря трём свойствам: полной интерпретируемости (каждое решение объяснимо), нулевому холодному старту (не требует размеченных данных) и детерминированности (одинаковый входной текст всегда получает одну и ту же метку)~\cite{Cohen1995}. Обзор~\cite{aggarwal2012survey} показывает, что для коротких текстов с фиксированным словарём (каковыми являются назначения платежей конкретного предприятия) rule-based подход с регулярными выражениями и словарём контрагентов обеспечивает точность, сравнимую с ML-классификаторами, но без затрат на разметку.

ML-подходы к классификации транзакций обоснованы в контексте персональных финансовых менеджеров (PFM): при наличии многомиллионных размеченных корпусов транзакций gradient boosting и нейросетевые классификаторы демонстрируют точность 90--95\%. Однако для отдельного МСП такие корпусы недоступны, а внутренний размеченный датасет содержит лишь несколько тысяч операций --- недостаточно для надёжного обучения без предобработки. В этом сценарии ML выгоден как второй уровень классификации поверх rule-based системы: операции, не покрытые правилами, передаются модели, а та, в свою очередь, принимает решение на основе признаков контрагента, суммы и времени~\cite{pustejovsky2012natural}.

Появление трансформерных архитектур~\cite{vaswani2017attention} и предобученных языковых моделей~\cite{devlin2018bert} поставило вопрос об их применимости к классификации банковских транзакций. Технически BERT и его производные могут быть дообучены (fine-tuned) на размеченном корпусе назначений платежей, однако у этого подхода есть ключевые ограничения. Во-первых, назначения платежей --- крайне короткие тексты с высокой долей числовых идентификаторов и аббревиатур, для которых токенизатор субтокенов работает неэффективно. Во-вторых, трансформеры не предоставляют детерминированных гарантий: два идентичных текста с разными контекстными фичами могут получить разные классы~\cite{ji2023survey}. В-третьих, inference трансформерной модели требует значительных вычислительных ресурсов, несовместимых с типичной инфраструктурой МСП.

Большие языковые модели (LLM) типа GPT-4~\cite{brown2020language} решают задачу классификации в режиме zero-shot или few-shot: без дообучения, через текстовый промпт с примерами. Для банковских транзакций это привлекательно на этапе первичной разметки, однако неприемлемо как производственный инструмент по причинам галлюцинаций~\cite{ji2023survey} (модель может уверенно возвращать несуществующую категорию), стохастичности (разные прогоны дают разные результаты) и конфиденциальности (API-запросы к внешним LLM передают финансовые данные третьей стороне). Локальное развёртывание LLM~\cite{brown2020language} частично снимает проблему приватности, но не гарантийность результата.

Таким образом, академический консенсус в данной области~\cite{aggarwal2012survey,sculley2015hidden,ji2023survey} сводится к следующему: для задачи классификации транзакций одного МСП оптимальна гибридная схема, где rule-based система закрывает 90--95\% операций, LLM применяется в качестве разметчика для редких случаев на этапе формирования датасета, а ML-классификатор становится актуальным по мере накопления размеченного корпуса.

\section{Прогнозирование спроса и денежных потоков}

Прогнозирование временны́х рядов --- одна из наиболее изученных областей машинного обучения. Фундаментальные подходы к прогнозированию систематизированы в~\cite{hyndman2021forecasting}: от ARIMA-семейства~\cite{box1970distribution} до нейросетевых архитектур. Для МСП особый интерес представляют результаты крупных прогнозных соревнований M4~\cite{makridakis2018m4} и M5~\cite{makridakis2022m5}, охватывающих десятки тысяч реальных бизнес-рядов.

Ключевой вывод M4~\cite{makridakis2018m4}: на горизонте прогнозирования до 12 месяцев гибридные модели (статистика + ML) превосходят чисто нейросетевые подходы. Более того, простые методы --- экспоненциальное сглаживание и скользящее среднее --- конкурентоспособны со сложными ансамблями при малом числе наблюдений. M5~\cite{makridakis2022m5} уточняет этот вывод для розничного прогнозирования: LightGBM~\cite{ke2017lightgbm} с правильно подобранными лаговыми признаками обходит статистические методы при объёме более 500 наблюдений на ряд, но проигрывает при меньшем объёме. Этот порог критически важен для МСП с 3--5 активными SKU и историей менее 18 месяцев.

Специфической проблемой является \textit{intermittent demand} (прерывистый спрос)~\cite{croston1972forecasting} --- ситуация, когда значительная доля наблюдений нулевая. Для хлебобулочной и кондитерской продукции это характерно на дневном уровне агрегации: возвраты поступают нерегулярно, отгрузки концентрированы в определённые дни недели. Метод Кростона~\cite{croston1972forecasting} разделяет прогноз ненулевого спроса и частоту его возникновения, однако смещён в большую сторону~\cite{syntetos2001bias}. Современная альтернатива --- hurdle-модели~\cite{lambert1992zero}, разделяющие задачу на бинарную классификацию (будет ли событие) и регрессию (каков объём), что позволяет использовать произвольные ML-методы на каждой стадии.

Прогнозирование денежного потока на уровне МСП --- значительно менее изученная задача, чем прогнозирование спроса. В отличие от спроса, денежный поток формируется пересечением нескольких процессов: производственных закупок, выручки с временным лагом, налоговых платежей по квартальному расписанию, нерегулярных операционных расходов. Применение методов машинного обучения к прогнозированию корпоративного cash flow в целом~\cite{lopezdeprado2018advances} основано на тех же принципах, что и прогнозирование спроса: формирование признакового пространства с лагами, скользящими статистиками и календарными переменными~\cite{ke2017lightgbm}. Принципиальное отличие --- высокая нерегулярность: для крупного предприятия платежи распределены относительно равномерно, для МСП один крупный платёж может составлять 30--40\% месячного оборота.

Анализ розничных временны́х рядов~\cite{fildes2019retail} показывает, что методы машинного обучения наиболее эффективны при наличии внешних ковариат --- сезонного индекса, праздничного календаря, данных о промоакциях. Для МСП с монозависимостью от одной торговой сети доля выручки от этого канала является сильнейшим предиктором месячного спроса.


\section{Проблемы учёта и ведения управленческой отчётности в МСП}

Помимо нехватки прогнозных и классификационных инструментов, малые и средние предприятия систематически испытывают более базовую проблему: само ведение учёта и сведение управленческой отчётности устроено фрагментарно. Данные о продажах, оплатах, возвратах и движении денежных средств хранятся в разных, слабо связанных друг с другом источниках --- учётной системе, банк-клиенте, ручных таблицах Excel разных сотрудников --- и для получения целостной картины собственнику или бухгалтеру приходится каждый раз вручную сводить их в единый отчёт.

Это создаёт две практические проблемы. Во-первых, отсутствует единое место, где управленческая отчётность собрана целиком и обновляется регулярно: подготовка сводки требует ручного труда, а значит выполняется редко (раз в месяц или квартал), и решения принимаются на основе устаревших цифр. Во-вторых, даже когда отчётность собрана, она представляет собой таблицы чисел без визуального и содержательного контекста --- трудно быстро увидеть динамику, аномалии или структурные сдвиги, не проводя отдельный анализ. Для МСП с ограниченным штатом (часто один бухгалтер или собственник, выполняющий эту функцию) это означает, что управленческие решения принимаются не вовремя и не на основе полной картины, а на основе интуиции и отдельных, не связанных между собой цифр.

Таким образом, прежде чем переходить к продвинутым задачам --- классификации транзакций, прогнозированию спроса и денежных потоков (рассмотренным в следующих разделах), --- необходимо решить более фундаментальную задачу: обеспечить единое, регулярно обновляемое и наглядное представление управленческой отчётности, на основе которого можно строить более сложную аналитику. Эта задача является сквозной для всей работы и непосредственно мотивирует создание собственной CRM-платформы, описанной в главе~\ref{ch:app}.

\section{Синтез обзора и позиционирование работы}

Обзор литературы позволяет выделить три принципиальных разрыва в существующих исследованиях. \textbf{Первый разрыв} --- между теоретическими возможностями ML-методов (исследуемыми на крупных корпоративных датасетах) и практическими ограничениями МСП (малый объём данных, отсутствие стандартизированных форматов, единственный аналитик). Большинство опубликованных работ исследуют ML на датасетах размером от 10\,000 наблюдений; реальная история МСП на уровне SKU~$\times$~месяц составляет 100--150 строк.

\textbf{Второй разрыв} --- между технической точностью моделей и управленческой интерпретируемостью результатов. Работы по финансовому ML~\cite{lopezdeprado2018advances,khandani2010consumer} оптимизируют метрики качества (AUC, WAPE, F1), оставляя за рамками вопрос о том, как переводить прогностический результат в конкретное управленческое действие (решение о закупке, размер резервного фонда, приоритет проверки операции).

\textbf{Третий разрыв} --- между академическими рекомендациями по MLOps~\cite{sculley2015hidden} и их практической реализацией в МСП. Воспроизводимые пайплайны, автоматический мониторинг дрейфа и версионирование моделей --- стандарты, сформированные на примере крупных технологических компаний с десятками инженеров; их адаптация к условиям одного аналитика требует принципиальных упрощений без потери сути.

Данная курсовая работа адресует все три разрыва: она исследует весь аналитический конвейер --- от ETL до прогнозных моделей --- на реальных данных конкретного МСП с типичными ограничениями объёма и качества данных, с явным акцентом на управленческую интерпретацию и минималистичный, но воспроизводимый MLOps-контур.

Структура работы напрямую отражает три выделенных разрыва. Первый разрыв (объём данных vs. требования ML-методов) разбирается через всю цепочку последующих глав: подготовка данных и классификация транзакций (главы~\ref{ch:data} и~\ref{ch:classification}) формируют максимально полный для конкретного МСП датасет, финансовый аудит (глава~\ref{ch:dds}) показывает, какие управленческие вопросы этот датасет должен закрывать, а обучение прогнозных моделей (глава~\ref{ch:ml}) явно сопоставляет сложные ML-алгоритмы с простыми бейзлайнами на данных малого объёма --- результат этого сопоставления (раздел об ограничениях главы~\ref{ch:ml}) эмпирически подтверждает наблюдение M4/M5~\cite{makridakis2018m4,makridakis2022m5} о конкурентоспособности простых методов при коротких рядах. Второй разрыв (точность модели vs. управленческая интерпретируемость) закрывается тем, что каждая модель в главе~\ref{ch:ml} сопровождается не только метрикой WAPE, но и явной управленческой интерпретацией (прогноз кассового разрыва, производственная рекомендация, приоритизированная очередь аномалий для бухгалтера). Третий разрыв (академический MLOps vs. ресурсы МСП) адресован практически: глава~\ref{ch:app} описывает не теоретическую архитектуру, а работающую самописную аналитическую систему, в которую интегрирован весь прогнозный контур --- от загрузки первичных данных до мониторинга денежного потока и транзакций, --- реализованную силами одного разработчика без выделенной инфраструктуры крупной компании.

Из обзора литературы следует методологически важный вывод, подтверждаемый далее на практике (главы~\ref{ch:dds}--\ref{ch:ml}): малый объём данных --- это ограничение на выбор метода, а не аргумент против внедрения прогнозных систем как таковых. Согласно M4 и M5~\cite{makridakis2018m4,makridakis2022m5}, при недостатке наблюдений именно простые и интерпретируемые модели (скользящие средние, hurdle-схемы, rule-based классификация) дают наибольшую устойчивость, и именно такие модели в первую очередь имеет смысл встраивать в производственный контур МСП. Иначе говоря, недостаточный для «промышленного» ML объём данных --- типичная характеристика любого малого предприятия, а не повод откладывать автоматизацию до накопления многотысячных датасетов: ценность для управления создаёт не абсолютная точность модели, а регулярность и воспроизводимость прогноза, заменяющего интуитивную оценку «как в прошлом периоде».
